\begin{enumerate}
\item Xác định các yếu tố bất biến của hình gốc: quan hệ song song, trung điểm, trọng tâm, tỉ số độ dài các đoạn thẳng cùng nằm trên một đường thẳng hoặc trên hai đường thẳng song song.
\item Chọn hình biểu diễn phù hợp theo quy tắc (hình vuông/chữ nhật/thoi $\to$ hình bình hành; tam giác $\to$ tam giác; hình thang $\to$ hình thang giữ nguyên tỉ số hai đáy).
\item Vẽ hình đảm bảo các yếu tố bất biến; dùng nét liền cho đường nhìn thấy và nét đứt cho đường bị che khuất.
\end{enumerate}
Lưu ý: Hình bình hành có thể biểu diễn hình vuông, hình chữ nhật, hình thoi nhưng không thể biểu diễn hình thang có hai đáy không bằng nhau.
